\begin{afterword}
\summaryhead

The post is set in an imaginary factory. Blue eggs, called ``bleggs,'' go in one bin and red cubes, called ``rubes,'' in another, and the two differ in other ways too: bleggs are furred, flexible and opaque. A purple egg with all the other blegg features is plainly a blegg, whatever the definition says. Asked why the bins exist at all, the senior sorter, Susan, at last reveals that bleggs contain vanadium and rubes palladium.

But vanadium is not a definition either. About 2 per cent of ordinary-looking bleggs contain palladium. Such an object goes in the rube bin, yet it will probably glow in the dark like a blegg. So ``Is this object a blegg?'' stands for different questions on different occasions, such as which bin the object goes in or whether it will glow, and if it stood for none, no one would care. The author applies this to arguments over whether atheism is a ``religion.'' They are really about empirical claims, and a dictionary cannot settle them, because redrawing a word's boundary does not move anything in the world.

\responsehead

The post makes one point and makes it clearly. A category is used to answer practical questions, and when the questions come apart, so does the category. The imaginary factory is built to show exactly this, which is what a made-up example is for, and the object that goes in the rube bin but glows like a blegg shows it well.

The point is older than the post, which does not claim otherwise. William James made it the rule of pragmatism in 1907, with a dispute about whether a man circling a tree goes round a squirrel: ``If no practical difference whatever can be traced, then the alternatives mean practically the same thing, and all dispute is idle.'' The post's ``If it weren't standing in for some query, you'd have no reason to care'' is the same test, applied to categories.

The application to real disputes is weaker than the factory example. The post says what people who call atheism a religion are ``really'' trying to argue, with ``(I think)'' and without quoting anyone. The argument is the hypothetical one from the author's earlier post ``Uncritical Supercriticality.'' The reader has to take the diagnosis on trust.

In short: a clear dramatization of the pragmatic test for disputes about categories, which James stated in 1907, applied to real arguments through a reading of opponents who are not quoted.
\end{afterword}
